\section{Instalación de Ventilación}

En esta unidad solo se recogen las particularidades de Ventilación. El uso general de la herramienta ya se ha explicado en los apartados anteriores, por lo que aquí se detallan únicamente las reglas, avisos y casos propios de esta instalación.

\subsection{Descripción general de Ventilación}

\begin{itemize}
\item Esta instalación sirve para dibujar y editar redes de ventilación dentro de Allplan.
\item Está pensada para trabajar con conductos y accesorios automáticos a partir de un recorrido polilineal.
\item Lo que la diferencia de otras instalaciones es la combinación de tipos de conducto, diámetros, ángulos permitidos y accesorios específicos de ventilación.
\end{itemize}

\subsection{Sistemas disponibles en Ventilación}

En la configuración actual de Ventilación aparecen estos sistemas:

\begin{itemize}
\item \texttt{Conducto Impulsion};
\item \texttt{Conducto Extraccion};
\item \texttt{Conducto Aislado};
\item \texttt{Conducto Recuperador}.
\end{itemize}

Las diferencias principales entre ellos son estas:

\begin{itemize}
\item \texttt{Conducto Impulsion} y \texttt{Conducto Extraccion} comparten geometría base, pero usan color y artículo distintos.
\item \texttt{Conducto Aislado} trabaja con conducto dinámico y con reducción como conexión característica.
\item \texttt{Conducto Recuperador} trabaja con otra lógica de generación y solo admite giros de \texttt{90 grados}.
\end{itemize}

Cuando el usuario cambia de sistema, la herramienta adapta automáticamente el comportamiento del trazado, el diámetro base y los accesorios disponibles.

\manualimage{CAPTURA-VENTILACION-01.png}{Paleta de Ventilación con los cuatro sistemas visibles}

\subsection{Distribución en Ventilación}

En la paleta existe el campo \texttt{Tipo de Distribucion}, pero en la configuración registrada de Ventilación no aparece una distribución funcional específica como sí ocurre en Agua.

Por eso, para esta instalación:

\begin{itemize}
\item \pendingtag\ debe validarse si este campo se usa realmente o si permanece sin aplicación práctica para el usuario.
\end{itemize}

\subsection{Diámetros y cambios de diámetro en Ventilación}

\subsubsection{Diámetros por sistema}

Los diámetros confirmados en la configuración actual son:

\begin{itemize}
\item \texttt{Conducto Impulsion}: \texttt{75 mm};
\item \texttt{Conducto Extraccion}: \texttt{75 mm};
\item \texttt{Conducto Aislado}: \texttt{150 mm} y \texttt{160 mm};
\item \texttt{Conducto Recuperador}: \texttt{160 mm}.
\end{itemize}

Esto significa que el diámetro disponible depende del sistema seleccionado.

\manualimage{CAPTURA-VENTILACION-02.png}{Selector de diámetro en Conducto Aislado mostrando 150 mm y 160 mm}

\subsubsection{Qué ocurre al modificar el diámetro}

El cambio de diámetro sigue el comportamiento general descrito en \hyperref[sec:cambiar-diametro]{Cambiar el diámetro}.

En Ventilación, después de aceptar el cambio, la herramienta recalcula el tramo y las uniones afectadas. Esto puede afectar no solo al tramo seleccionado, sino también a las reducciones, conexiones, manguitos o codos que dependan de ese diámetro.

Cuando el cambio de diámetro se produce entre dos tramos consecutivos en línea recta:

\begin{itemize}
\item en conductos compatibles, la herramienta genera automáticamente la transición o la unión correspondiente.
\end{itemize}

En los sistemas donde existen accesorios automáticos, conviene revisar en pantalla cómo se rehace la unión después del cambio.

\subsubsection{Caso particular de Conducto Aislado}

En \texttt{Conducto Aislado}, la conexión característica registrada es la \texttt{reducción}.

Por eso, cuando el usuario trabaja con \texttt{150 mm} y \texttt{160 mm}, conviene revisar explícitamente la transición resultante.

\subsection{Accesorios automáticos en Ventilación}

Los accesorios automáticos dependen del sistema seleccionado.

En la configuración actual se confirma este comportamiento:

\begin{itemize}
\item \texttt{Conducto Impulsion} y \texttt{Conducto Extraccion}: \texttt{manguito} y \texttt{difusor};
\item \texttt{Conducto Aislado}: \texttt{reducción};
\item \texttt{Conducto Recuperador}: \texttt{conexion} y \texttt{codo 90}.
\end{itemize}

Los accesorios registrados para esta instalación son:

\begin{itemize}
\item \texttt{Manguito};
\item \texttt{Difusor};
\item \texttt{Conexion};
\item \texttt{Codo 90}.
\end{itemize}

\subsubsection{Manguitos}

En los sistemas normales, el manguito no actúa solo como accesorio visual.

También actúa como frontera de grupo para numeración y agrupación interna de la instalación.

\manualimage{CAPTURA-VENTILACION-03.png}{Ejemplo de Conducto Impulsion con manguito automático}

\subsubsection{Difusores}

En los sistemas que los usan, los difusores forman parte del comportamiento automático previsto por la instalación.

Conviene revisar su colocación final cuando el trazado cambia de dirección o de plano.

\manualimage{CAPTURA-VENTILACION-04.png}{Ejemplo de Conducto Extraccion con difusor}

\subsubsection{Recuperador con conexión y codo}

En \texttt{Conducto Recuperador}, el flujo de creación trabaja con un conjunto de piezas encadenadas:

\begin{itemize}
\item conducto;
\item conexión;
\item codo.
\end{itemize}

Esto lo diferencia claramente del resto de sistemas de Ventilación.

\manualimage{CAPTURA-VENTILACION-05.png}{Recuperador con conexión y codo 90}

\subsubsection{Limitaciones o incompatibilidades}

\texttt{Conducto Recuperador} no sigue la misma lógica de ángulos que el resto.

\begin{itemize}
\item \pendingtag\ debe validarse si aparecen avisos específicos cuando el usuario intenta forzar una combinación no compatible.
\end{itemize}

\subsection{Reglas geométricas y validaciones en Ventilación}

\subsubsection{Ángulos permitidos}

Los ángulos generales soportados en la instalación son:

\begin{itemize}
\item \texttt{45 grados};
\item \texttt{90 grados}.
\end{itemize}

Por sistema, la regla actual es esta:

\begin{itemize}
\item \texttt{Conducto Impulsion} y \texttt{Conducto Extraccion}: \texttt{45} y \texttt{90 grados};
\item \texttt{Conducto Aislado}: \texttt{45} y \texttt{90 grados};
\item \texttt{Conducto Recuperador}: solo \texttt{90 grados}.
\end{itemize}

\subsubsection{Longitud mínima}

La longitud mínima por sistema es:

\begin{itemize}
\item \texttt{Conducto Impulsion} y \texttt{Conducto Extraccion}: \texttt{300 mm};
\item \texttt{Conducto Aislado}: \texttt{300 mm};
\item \texttt{Conducto Recuperador}: \texttt{350 mm}.
\end{itemize}

\pendingtag\ El mensaje exacto mostrado cuando no se cumple la longitud mínima debe validarse en Allplan.

\subsubsection{Longitud máxima}

Ventilación no tiene una longitud máxima de tramo definida.

A diferencia de Agua, la herramienta no divide automáticamente un conducto largo por superar un límite máximo. En esta instalación, la validación de longitud documentada es la longitud mínima de cada sistema.

\subsubsection{Orientación 3D y cambios de plano}

En Ventilación, la orientación 3D del tramo afecta al resultado de:

\begin{itemize}
\item conductos;
\item codos;
\item conexiones;
\item difusores.
\end{itemize}

Además, existen casos específicos para:

\begin{itemize}
\item segmentos verticales;
\item y cambios de plano.
\end{itemize}

Por eso, no basta con revisar la planta. En esta instalación conviene comprobar también cómo se resuelve la orientación final del elemento.

\manualimage{CAPTURA-VENTILACION-08.png}{Orientación 3D o tramo vertical/inclinado}

\subsection{Layers y atributos específicos de Ventilación}

\subsubsection{Layers base}

Los layers base confirmados en el registro de Ventilación son:

\begin{itemize}
\item \texttt{IS\_CON\_VENT\_FAB};
\item \texttt{IS\_CON\_VENT\_FAB\_SOB1};
\item \texttt{IS\_CON\_VENT\_OBRA};
\item \texttt{IS\_CON\_VENT\_EIX};
\item \texttt{IS\_NOM\_CONDUCTE\_VENTILACIO}.
\end{itemize}

Esto cubre tanto los conductos como la capa de eje y la capa de etiqueta o nombre del conducto.

\manualimage{CAPTURA-VENTILACION-06.png}{Ejemplo de layer aplicado en Ventilación}

\subsubsection{Atributos detectados en scripts}

En los scripts de Ventilación se detectan, entre otros, estos atributos:

\begin{itemize}
\item \texttt{pmp\_pare};
\item \texttt{6\_CC\_IS};
\item \texttt{pmp\_altura};
\item \texttt{pmp\_amplada};
\item \texttt{pmp\_area};
\item \texttt{pmp\_diametre};
\item \texttt{pmp\_longitud};
\item \texttt{pmp\_nom};
\item \texttt{pmp\_seccio};
\item y otros atributos \texttt{pmp\_*} de clasificación y propiedades físicas.
\end{itemize}

\subsubsection{Casos concretos observados}

En la configuración actual se observan estos casos:

\begin{itemize}
\item \texttt{conducto\_normal} usa \texttt{6\_CC\_IS = IS08};
\item \texttt{conducto\_recuperador} usa \texttt{6\_CC\_IS = IS};
\item varias piezas inicializan \texttt{pmp\_pare} vacío.
\end{itemize}

Esto conviene tenerlo presente si el proyecto organiza información por atributo padre o por atributos de clasificación.

\subsection{Elementos definidos en Ventilación}

Aquí hay una inconsistencia que conviene documentar con claridad.

Por un lado, el registro de Ventilación declara como elementos 3D propios:

\begin{itemize}
\item \texttt{Manguito};
\item \texttt{Difusor};
\item \texttt{Conducto recuperador};
\item \texttt{Conexion};
\item \texttt{Codo 90}.
\end{itemize}

Pero, por otro lado, la página de \texttt{Elemento} del archivo \texttt{ventilacion\_polyline.pyp} muestra actualmente solo:

\begin{itemize}
\item \texttt{Caixa Connexions 200}.
\end{itemize}

Esto indica una probable herencia o copia de configuración desde otra instalación.

Por tanto:

\begin{itemize}
\item \pendingtag\ debe validarse qué elemento definido debe ver realmente el usuario en Ventilación;
\item \pendingtag\ conviene revisar si la paleta de Elemento en Ventilación debe corregirse antes de cerrar el manual.
\end{itemize}

\manualimage{CAPTURA-VENTILACION-10.png}{Página de Elemento mostrando el estado real en Ventilación}

\subsection{Macros en Ventilación}

La instalación incluye la misma lógica general de macros que otras instalaciones:

\begin{itemize}
\item selección de punto;
\item selección de macro SmartSymbol;
\item cota manual o relativa a local;
\item confirmación de inserción.
\end{itemize}

En la paleta de Ventilación se ven estos campos relacionados:

\begin{itemize}
\item selección de punto;
\item selección de local;
\item cota Z o altura sobre piso;
\item selección de macro.
\end{itemize}

\pendingtag\ si en Ventilación existe un uso específico de macros distinto del comportamiento compartido.

\manualimage{CAPTURA-VENTILACION-09.png}{Página de macros de Ventilación}

\subsection{Soportes en Ventilación}

Ventilación dispone de página específica de soportes.

Los campos confirmados en paleta son:

\begin{itemize}
\item \texttt{Tipo de soporte};
\item \texttt{Subtipo instalación};
\item \texttt{Superficie};
\item \texttt{Cota A};
\item \texttt{Cota B};
\item \texttt{Ángulo de inclinación};
\item modos de edición de soportes;
\item y \texttt{Atributo de soporte}.
\end{itemize}

Según el ejemplo de configuración disponible:

\begin{itemize}
\item existe al menos un caso de soporte \texttt{Omega};
\item el subtipo asociado puede ser \texttt{Ventilacion};
\item la superficie puede ser \texttt{Perforado}.
\end{itemize}

Las reglas prácticas de uso son estas:

\begin{itemize}
\item conviene definir tipo, superficie y cotas antes de insertar;
\item los soportes pueden acumularse y crearse después;
\item el modo \texttt{Edición Mover} permite recolocar soportes ya acumulados o seleccionados.
\end{itemize}

\pendingtag\ el listado completo de tipos de soporte realmente habilitados en Ventilación.

\manualimage{CAPTURA-VENTILACION-07.png}{Bloque de soportes con subtipo Ventilacion}

\subsection{Copias, duplicados y comportamiento especial en Ventilación}

Ventilación usa la misma lógica general de copias y agrupación que el resto de instalaciones basadas en \texttt{PolyLib}.

Los comportamientos relevantes observados son estos:

\begin{itemize}
\item los elementos pueden agruparse por tramos y por fronteras de manguito;
\item la numeración cambia cuando aparece un nuevo grupo;
\item la serialización del estado permite restaurar la instalación en edición.
\end{itemize}

El usuario debe entender especialmente esto:

\begin{itemize}
\item un manguito no es solo un accesorio visual;
\item también afecta a la agrupación lógica de la instalación.
\end{itemize}

Si el proyecto usa atributo padre y copia a otros archivos, esa lógica sigue siendo la global del sistema.

\pendingtag\ si Ventilación está usando en producción copias automáticas a otros archivos de dibujo y en qué casos.

\subsection{Mensajes y avisos propios de Ventilación}

A nivel de código se detectan mensajes técnicos y trazas de depuración, pero todavía no queda claro qué avisos funcionales ve el usuario final en Allplan.

Por eso, en esta instalación conviene validar en entorno real:

\begin{itemize}
\item el aviso por longitud mínima;
\item el aviso por selección o inserción inválida;
\item el aviso por combinación de geometría no soportada en recuperador;
\item y los mensajes de finalización si falta configuración de layers o atributos.
\end{itemize}

\subsection{Puntos a revisar con especial atención en Ventilación}

Antes de finalizar una instalación de Ventilación conviene revisar especialmente:

\begin{itemize}
\item que el sistema activo sea realmente \texttt{Conducto Impulsion}, \texttt{Conducto Extraccion}, \texttt{Conducto Aislado} o \texttt{Conducto Recuperador};
\item que el diámetro elegido sea compatible con ese sistema;
\item que en \texttt{Conducto Aislado} la transición entre \texttt{150 mm} y \texttt{160 mm} se haya resuelto correctamente;
\item que en \texttt{Conducto Recuperador} los giros respondan a la lógica de \texttt{90 grados};
\item que los layers aplicados correspondan al criterio del plano;
\item que los atributos de clasificación se hayan aplicado como se espera;
\item que la página de \texttt{Elemento} esté mostrando el elemento correcto para esta instalación;
\item y que los soportes se hayan revisado con subtipo, superficie, cotas y tipo correctos.
\end{itemize}

